\subsection{DNS重绑定攻击}

DNS重绑定攻击(DNS Rebinding)是一种利用域名解析机制发起的网络攻击方法，主要利用域名解析的应答响应报文，规避浏览器的同源防御策略，
获取受害者内网资源。

DNS重绑定攻击的实用场景未受害者处于安全受控的内网环境，攻击者难以直接从外部网络环境对受害者内网资源进行访问。
具体攻击流程为：
\begin{itemize}
    \item 攻击者引导受害者访问其受控制的恶意域名，恶意域名服务器将该域名指向一个公开可访问的主机地址，并将恶意域名资源记录的有效TTL设置为极短。
    \item 恶意网站中包含JS的恶意代码，但是会受到同源策略的限制。
    \item 当资源记录过期时，攻击者控制的恶意域名服务器将该域名解析为受害者所在的内部网络的主机地址，由于此时浏览器认为整个网络访问过程属于同一个域名（这就绕过了浏览器同源策略，其中同指协议、端口、主机都相同），
    将允许恶意JS代码访问受害者内部网络环境。
\end{itemize}